\begin{abstract}
Selective labels can hide the losses that determine a predictor's population tail risk. Even unlimited historical data can leave its conditional value-at-risk (CVaR) unresolved. For binary outcomes and bounded losses, we derive the sharp population CVaR interval under an outcome-dependent odds-ratio sensitivity model: every value is compatible with the observations and the stated restriction. We construct calibrated statistical intervals enclosing every compatible risk simultaneously across candidate predictors at fixed sensitivity and tail levels, using cross-fitted bias correction. These certificates separate persistent hidden-outcome ambiguity from sampling error and support risk limits, predictor comparisons, and conservative selection. We also derive the minimum-cost revelation allocation for a fixed predictor and loss threshold. Tail identification value measures ambiguity removed when audited contributions become known and the original sensitivity bounds apply separately to audited and still-hidden outcomes. Synthetic and lending experiments examine endpoint accuracy, population assessment, and targeted audits. Under controlled masking of known repayment outcomes, revealing about 16\% of hidden labels removes a median 61\% of estimated fixed-threshold ambiguity, versus 16\% for random revelation. The framework specifies which tail-risk claims the evidence supports and where additional labels improve precision.
\end{abstract}
